\section{Common Crawl：web-scale raw source}

前面的法律与治理字段说明“能抓到”不等于“可直接训练”，本节用 Common Crawl 展开 web-scale raw source 的技术链。Crawler 从 seed URL 维护 frontier、下载页面并发现新链接；WARC/WET、HTML extraction、语言识别和去重随后逐层改变数据分布。

\begin{figure}[H]
\centering
\includegraphics[width=0.5\textwidth]{webcrawler-architecture.png}
\caption{Web crawler 架构：从 seed URLs 出发，下载页面并把超链接加入队列。}
\figsource{Wikipedia WebCrawlerArchitecture，本地化后供 CS336 Lecture 13 使用。}
\end{figure}

\begin{knowledgebox}{课堂提示：Common Crawl 是原料，不是数据集答案}
课上提醒，直接说“我们用 Common Crawl”几乎没有复现价值。至少要报告 snapshot、URL/domain policy、WARC-to-text 工具、过滤版本、去重粒度和最终 token 数；同一个 crawl 经不同 extraction 与 filtering 可以得到行为明显不同的模型。
\end{knowledgebox}

\begin{knowledgebox}{读图：crawler 不是简单 wget 全网}
Crawler 需要 selection policy、politeness policy、re-visit policy。它必须决定爬什么、多久重爬、是否尊重 robots.txt、如何避免压垮网站。Common Crawl 每月增加数十亿网页，但 crawl 本身仍只是 raw source。
\end{knowledgebox}

Common Crawl 提供 WARC 和 WET。WARC 保存 raw HTTP response，例如 HTML；WET 是转换成 text 后的结果。HTML-to-text 是 lossy process，会影响最终模型。

\begin{figure}[H]
\centering
\includegraphics[width=0.55\textwidth]{dclm-wet.png}
\caption{HTML-to-text 工具选择会影响下游 LM accuracy。}
\figsource{CS336 Lecture 13 官方源引用图。}
\end{figure}

\begin{knowledgebox}{读图：为什么 WET 不是无害中间格式}
HTML 包含正文、导航、广告、脚本、评论、页脚。不同 extraction 工具会保留或丢弃不同内容。若正文抽取差，模型会学到网页模板噪声；若抽取过严，又会丢掉有用长尾内容。
\end{knowledgebox}

\begin{knowledgebox}{术语消化：Common Crawl 到训练文本的处理链}
\begin{tabular}{L{0.22\textwidth}L{0.34\textwidth}L{0.34\textwidth}}
\toprule
步骤 & 做什么 & 主要风险 \\
\midrule
HTML extraction & 从 WARC/HTML 中抽正文。 & boilerplate、广告、导航、评论混入或正文丢失。 \\
Language ID & 判断文档语言。 & 多语混合、代码混合、低资源语言误判。 \\
Quality filtering & 用规则或模型判断是否像高质量文本。 & 过滤掉长尾有用内容，引入 Wikipedia/指令数据偏差。 \\
Deduplication & 移除重复或近重复文档/段落。 & 太弱会重复训练，太强会误删模板相似但内容不同的文本。 \\
Toxicity/PII filtering & 去除有害内容或个人信息。 & 过强会损害分布覆盖，过弱会带来安全和隐私风险。 \\
\bottomrule
\end{tabular}
\end{knowledgebox}

